% !PRINTED 174
\weberpage{174}

%<<CONT
und ergeben also nach (16) des vorigen Paragraphen
\[
\alpha s_0 + \beta s_1 + \gamma s_2 = 0,
\]
\[
\alpha^2 s_0 + 2\alpha\beta s_1 + 2\alpha\gamma s_2 + \beta^2 s_2 + 2\beta\gamma s_3 + \gamma^2 s_4 = 0,
\]
worin $s_0 = 3$ ist. Multiplicirt man die letztere Gleichung mit $s_0$ und zieht das Quadrat der ersten davon ab, so folgt
\weq{6}{
\beta^2(s_0 s_2 - s_1^2) + 2\beta\gamma(s_0 s_3 - s_1 s_2) + \gamma^2(s_0 s_4 - s_2^2) = 0.
}
Die Discriminante dieser quadratischen Gleichung
\[
(s_0 s_3 - s_1 s_2)^2 - (s_0 s_2 - s_1^2)(s_0 s_4 - s_2^2)
\]
giebt entwickelt
\[
-s_0(s_0 s_2 s_4 - s_0 s_3^2 - s_1^2 s_4 - s_2^3 + 2 s_1 s_2 s_3),
\]
ist also, wenn $D$ die Discriminante der cubischen Gleichung (1) bedeutet [\S.\,46, (9)], gleich $-3D$, und die Auflösung von (6) giebt also
\weq{7}{
\frac{\beta}{\gamma} = \frac{-(s_0 s_3 - s_1 s_2) + \sqrt{-3D}}{s_0 s_2 - s_1^2}.
}
Damit ist die Gleichung (3) auf eine reine Gleichung zurückgeführt, die zu ihrer Lösung nur noch das Ziehen einer Cubikwurzel erfordert. Welches Vorzeichen wir der in (7) vorkommenden Quadratwurzel geben wollen, steht in unserem Belieben.

Um die \gesperrt{Tschirnhausen}-Transformation für die biquadratische Gleichung zu benutzen, kann man etwa so verfahren, dass man in der umgeformten Gleichung
\weq{8}{
y^4 + b_1 y^3 + b_2 y^2 + b_3 y + b_4 = 0
}
die $\alpha_0, \alpha_1, \alpha_2, \alpha_3$ so bestimmt, dass
\weq{9}{
b_1 = 0, \quad b_3 = 0
}
wird, wodurch (8) in eine quadratische Gleichung für $y^2$ übergeht, aus deren Wurzeln noch die Quadratwurzeln gezogen werden müssen, so dass, nachdem die Gleichungen (9) gelöst sind, noch zweimal eine Quadratwurzel gezogen werden muss.

Wenn man mittelst der ersten Gleichung (9) aus der zweiten $\alpha_0$ eliminirt, so entsteht eine homogene cubische Gleichung zwischen $\alpha_1, \alpha_2, \alpha_3$. Man kann daher eines von den beiden Verhältnissen $\alpha_1:\alpha_2:\alpha_3$ beliebig annehmen, z. B. $\alpha_3 = 0$ setzen, und erhält zur Bestimmung des anderen eine \gesperrt{cubische Gleichung}. Dadurch ist dann die Lösung der biquadratischen Gleichung auf die einer cubischen und auf Quadratwurzeln zurückgeführt.
